\chapter{Random functional exposures and the covariance envelope}
\label{chap:random-exposure}

This is the trunk of the development.
An asset's factor loading is a probability law rather than a coefficient, and
the quadratic transport distance between two such laws is exactly the geometry
that polarization converts into a covariance bound.
Everything in the later chapters is either a specialization of this chapter or a
separate model that borrows its polarization step.

\section{The exposure law and its risk coordinates}

\begin{definition}[Random functional exposure law]
  \label{def:exposure-law}
  Asset \(i\) carries a random loading \(B_i\sim P_i\in\mathcal P_2(\mathcal H)\)
  in a real Hilbert space \(\mathcal H\), and the common factor process has
  covariance operator \(\Gamma\).
  The risk coordinates and factor-weighted second moments are
  \[
    Z_i=\Gamma^{1/2}B_i,\qquad v_i=\mathbb E\lVert Z_i\rVert^2 .
  \]
  For a coupling \(\pi\in\Pi(P_i,P_j)\) the systematic covariance is
  \(\kappa^\pi_{ij}=\mathbb E_\pi\langle Z_i,Z_j\rangle\), and the
  factor-weighted quadratic transport cost is
  \[
    W_{2,\Gamma}^2(P_i,P_j)
      =\inf_{\pi\in\Pi(P_i,P_j)}\mathbb E_\pi\lVert Z_i-Z_j\rVert^2 .
  \]
  Every object above is now formalized at this level, not only interpreted:
  risk coordinates as a pushforward under \(\Gamma^{1/2}\), the second moment
  and systematic covariance as Bochner integrals, and the transport cost as the
  infimum defining \(W_2\) on \(\mathcal P_2(\mathcal H)\).
  The finite instance of \Cref{def:finite-exposure-cloud} is the estimator of
  this object rather than a separate model, and the two agree exactly on
  finitely supported laws.

  One caveat survives and is unchanged: existence of an optimal plan on
  \(\mathcal P_2(\mathcal H)\) is \emph{not} proved.
  It is supplied as a certificate wherever an optimum is needed, so every
  attainment statement below reads ``given such a plan'', never ``one exists''.
\end{definition}

\begin{definition}[Finite exposure cloud and coupling]
  \label{def:finite-exposure-cloud}
  \uses{def:exposure-law}
  A finite exposure cloud is a finite support together with nonnegative masses
  summing to one, and a coupling of two clouds is a nonnegative mass matrix with
  the two clouds as marginals.
  Over a coupling this fixes the systematic covariance
  \(\kappa^\pi_{ij}\), the transport cost \(C(\pi)\), the second moments
  \(v_i\), and the barycentre, with no measure-theoretic optimal-transport
  input.
  A finite optimal-cost certificate is a separate datum: a coupling together
  with the assertion that its cost is least in the coupling set.
\end{definition}

\begin{lemma}[Polarization]
  \label{lem:polarization}
  \ppClaimBinding{lem:polarization}
  For any \(x,y\) in a real inner-product space,
  \[
    \langle x,y\rangle
      =\tfrac12\left(\lVert x\rVert^2+\lVert y\rVert^2-\lVert x-y\rVert^2\right).
  \]
\end{lemma}

\begin{proof}
  \ppClaimProofBinding{lem:polarization}
  Expand the squared norm of the difference.
\end{proof}

\begin{lemma}[Transport polarization for a coupling]
  \label{lem:transport-polarization}
  \uses{def:finite-exposure-cloud}
  For every coupling \(\pi\) of two finite exposure clouds,
  \[
    \kappa^\pi_{ij}=\tfrac12\left(v_i+v_j-C(\pi)\right).
  \]
  The identity is exact and holds coupling by coupling; no optimality is
  involved.
\end{lemma}

\begin{proof}
  \uses{lem:polarization}
  Apply \Cref{lem:polarization} inside the coupling expectation and use that
  the marginals reproduce \(v_i\) and \(v_j\).
\end{proof}

\section{The Fréchet class and its endpoints}

\begin{definition}[Fréchet class of attainable covariances]
  \label{def:frechet-class}
  \uses{def:finite-exposure-cloud}
  The Fréchet class of a pair of clouds is the set of systematic covariances
  \(\kappa^\pi_{ij}\) attained as \(\pi\) ranges over their couplings.
  The coupling set is convex and \(\kappa^\pi_{ij}\) is affine in \(\pi\), so
  the class is an interval whenever both endpoints are attained.
\end{definition}

\begin{theorem}[Covariance envelope from optimal transport]
  \label{thm:random-exposure-envelope}
  \ppClaimBinding{thm:random-exposure-envelope}
  \uses{def:frechet-class}
  Given a finite optimal transport-cost certificate, the ceiling
  \[
    \kappa^{\max}_{ij}=\tfrac12\left(v_i+v_j-W_{2,\Gamma}^2(P_i,P_j)\right)
  \]
  is the greatest systematic covariance attainable over the coupling set.
  This is a statement about what the two marginal laws permit; it does not
  identify the realized covariance, which requires the economic joint law.
  It is the \emph{upper} endpoint only; the lower one is
  \Cref{thm:random-exposure-lower-endpoint}, a separate theorem.
\end{theorem}

\begin{proof}
  \ppClaimProofBinding{thm:random-exposure-envelope}
  \uses{lem:transport-polarization}
  \Cref{lem:transport-polarization} turns least cost into greatest covariance
  monotonically, and the certificate supplies leastness.
\end{proof}

\begin{theorem}[Reflected lower endpoint]
  \label{thm:random-exposure-lower-endpoint}
  \ppClaimBinding{thm:random-exposure-lower-endpoint}
  \uses{def:frechet-class}
  Applying the same argument to the reflected law \((-I)_\#P_j\) gives the
  sharp lower endpoint of the Fréchet class.
\end{theorem}

\begin{proof}
  \ppClaimProofBinding{thm:random-exposure-lower-endpoint}
  \uses{lem:transport-polarization}
  Covariance is odd and the second moment even under reflection of one
  argument, so the greatest covariance against the reflected law is the least
  covariance against the original.
\end{proof}

\begin{theorem}[Transport-excess gap]
  \label{eq:random-exposure-gap}
  \ppClaimBinding{eq:random-exposure-gap}
  \uses{def:finite-exposure-cloud}
  For any coupling,
  \(\kappa^\pi_{ij}=\kappa^{\max}_{ij}-\tfrac12\operatorname{excess}(\pi)\).
  The checked statement is an identity in the reference cost: it holds for any
  scalar placed where the optimal cost sits, and therefore asserts nothing
  about the sign of the excess on its own.
\end{theorem}

\begin{proof}
  \ppClaimProofBinding{eq:random-exposure-gap}
  \uses{thm:random-exposure-envelope, lem:transport-polarization}
  Subtract the reference cost from the coupling cost inside
  \Cref{lem:transport-polarization} and rearrange.
\end{proof}

\begin{theorem}[The excess is nonnegative]
  \label{eq:random-exposure-excess-nonneg}
  \ppClaimBinding{eq:random-exposure-excess-nonneg}
  \uses{def:frechet-class}
  Against a certified optimal transport cost, \(\operatorname{excess}(\pi)\ge0\)
  for every coupling.
  This is what makes the gap a one-sided moment inequality and hence testable;
  it is a separate theorem from the identity above, because the identity holds
  for any reference cost while the sign needs that cost to be the infimum.
\end{theorem}

\begin{proof}
  \ppClaimProofBinding{eq:random-exposure-excess-nonneg}
  \uses{def:frechet-class}
  Immediate from minimality of the certified cost over the coupling set.
\end{proof}

\begin{proposition}[Pairwise envelopes do not assemble]
  \label{prop:pairwise-not-psd}
  \notready
  \uses{thm:random-exposure-envelope, thm:joint-coherence}
  The matrix of independently optimal pairwise envelopes
  \(\left[\kappa^{\max}_{ij}\right]\) need not be positive semidefinite,
  because pairwise-optimal couplings may be mutually inconsistent.
  A globally coherent covariance requires one joint law, as in
  \Cref{thm:joint-coherence}.
  Stating this boundary is part of the contribution and is not a defect of the
  envelope.
\end{proposition}

\section{Transfer from observable characteristics}

\begin{definition}[Randomized bi-Lipschitz pushforward]
  \label{ass:random-common-map}
  \uses{def:exposure-law}
  The exposure law is generated from an observable characteristic law by one
  common stochastic map with bounded slack:
  \[
    B_i=T(X_i,U_i)=t(X_i)+\delta_i(X_i,U_i),\qquad X_i\sim C_i,\qquad
    \lVert\delta_i(X_i,U_i)\rVert\le\tau_i,
  \]
  with the common carrier \(t\) bi-Lipschitz,
  \(\ell\,d(x,y)\le\lVert t(x)-t(y)\rVert\le L\,d(x,y)\) for
  \(0<\ell\le L\).
  The constants \(L\), \(\ell\), and the slack radii \(\tau_i\) are maintained
  modelling restrictions, not quantities the observed distances identify or the
  empirical design estimates.
\end{definition}

\begin{theorem}[Upper transport transfer under slack]
  \label{eq:random-w2-upper-transfer}
  \ppClaimBinding{eq:random-w2-upper-transfer}
  \uses{ass:random-common-map, def:finite-exposure-cloud}
  Under the upper Lipschitz bound and bounded per-endpoint slack, the optimized
  squared exposure transport cost of two finitely supported clouds is at most
  \(L^2\) times the optimized squared characteristic cost, plus a slack term
  carrying the maintained support diameter.
\end{theorem}

\begin{proof}
  \ppClaimProofBinding{eq:random-w2-upper-transfer}
  \uses{lem:transport-polarization}
  Push the certificate coupling of the characteristic clouds through \(T\),
  bound the resulting cost pointwise, and absorb the slack endpoints.
\end{proof}

\begin{theorem}[Lower transport transfer under slack]
  \label{eq:random-w2-lower-transfer}
  \ppClaimBinding{eq:random-w2-lower-transfer}
  \uses{ass:random-common-map, def:finite-exposure-cloud}
  Under a positive lower Lipschitz constant and bounded per-endpoint slack, the
  optimized squared exposure transport cost is at least \(\ell^{-2}\) times the
  optimized squared characteristic cost, less a slack term carrying the
  maintained support diameter.
\end{theorem}

\begin{proof}
  \ppClaimProofBinding{eq:random-w2-lower-transfer}
  \uses{lem:transport-polarization}
  Symmetric to the upper direction, using injectivity of \(T\) supplied by the
  lower bound.
\end{proof}

\begin{corollary}[Metric-level bracket]
  \label{cor:random-w2-metric}
  \notready
  \uses{eq:random-w2-upper-transfer, eq:random-w2-lower-transfer}
  At the level of the metric rather than the squared optimized cost, the slack
  enters additively and no support diameter appears:
  \[
    \ell^{-1}W_2(C_i,C_j)-(\tau_i+\tau_j)
      \le W_{2,\Gamma}(P_i,P_j)
      \le L\,W_2(C_i,C_j)+(\tau_i+\tau_j).
  \]
  This sharper form is what the empirical diagnostic uses.

  The upper direction is machine-checked directly at this level: two
  synchronous-coupling legs bound the slack, one contraction leg carries the
  transmitted laws, and the \(\mathcal P_2(\mathcal H)\) triangle inequality
  chains them.
  The lower direction still needs the reverse Lipschitz half, which supplies
  injectivity of \(T_\Gamma\) on its image, and remains a prose result.
  The node keeps \texttt{notready} because it asserts both directions.
\end{corollary}

\begin{corollary}[Covariance bracket the diagnostic tests]
  \label{eq:random-kappa-bracket}
  \notready
  \uses{thm:random-exposure-envelope, cor:random-w2-metric}
  Substituting the metric bracket into the envelope brackets
  \(\kappa^{\max}_{ij}\) in observable characteristic distances and the
  maintained \((L,\ell,\tau)\).
  Coverage of this bracket, and the share of unresolved dyads, is the primary
  reported diagnostic.
\end{corollary}

\section{Coherent aggregation}

\begin{theorem}[One joint law gives a coherent covariance]
  \label{thm:joint-coherence}
  \uses{def:finite-exposure-cloud}
  A joint exposure law over all assets reproduces each pairwise systematic
  covariance on its diagonal coupling, and the induced portfolio quadratic form
  equals the second moment of the aggregated loading
  \(B_w=\sum_i w_iB_i\).
  Coherence is therefore automatic once the object is a single joint law, and
  the quadratic form is positive semidefinite for that reason alone.
  This is the honest route around \Cref{prop:pairwise-not-psd}.
\end{theorem}

\begin{proof}
  \uses{lem:polarization}
  Expand the quadratic form of the aggregate and identify terms with the
  marginal and pairwise objects.
\end{proof}

\begin{theorem}[Factor-return semantics of the covariance]
  \label{thm:factor-return-semantics}
  \uses{thm:joint-coherence}
  With a centred isotropic factor and residuals cross-orthogonal to the
  exposures, the covariance of factor-driven returns equals the systematic
  covariance of the exposure laws, and the residual cross term vanishes.
  This is what licenses reading \(\kappa_{ij}\) as a return covariance rather
  than only as a loading-space inner product; without the orthogonality
  hypothesis the whole-return claim does not follow.
\end{theorem}

\begin{proof}
  \uses{thm:joint-coherence}
  Expand the return cross moment and kill the residual terms with the
  orthogonality hypothesis.
\end{proof}

\section{Rehomed from the retired deterministic chapter}

The retirement recorded on 2026-08-17 withdrew the manifest entries for the
deterministic branch and kept exactly two, because the pivot keeps their claims.
Those two moved here on 2026-08-21 when that chapter was deleted.
Both are Paper~1 claims, so the trunk is where they belong; neither statement
nor declaration changed in the move.

\begin{corollary}[Approximate common map]
  \label{cor:approximate-map-bound}
  \uses{ass:random-common-map}
  If each exposure deviates from a common map value by at most \(\tau_i\), the
  covariance floor holds with \(LD\) inflated to \(LD+\tau_i+\tau_j\), and exact
  transmission is the \(\tau=0\) case.

  The statement fixes no characteristic distance: \(D\) is an arbitrary
  nonnegative dissimilarity, and energy distance was only ever one instance.
  Hosting it under the energy development misread that generality, which is why
  the declaration now lives in \texttt{Transmission.Floor} and the node here.
  Contrast \Cref{thm:random-exposure-envelope}, which is \emph{not}
  distance-agnostic: the sharp envelope needs quadratic transport specifically,
  because minimising the coupling cost is the same optimisation polarisation
  turns into an inner product.
\end{corollary}

\begin{proof}
  \uses{lem:polarization}
  Triangle-inequality the deviations into the transmission bound, then polarize.
\end{proof}

\begin{theorem}[Portfolio variance floor]
  \label{eq:portfolio-envelope}
  \uses{cor:approximate-map-bound}
  For a finite book with nonnegative weights, the pairwise floors aggregate
  termwise into a floor on the squared systematic risk.
  Nonnegativity is a hypothesis, not a convention: signed weights can reverse
  individual pairwise inequalities and there is no signed-weight counterpart.
  No weight normalization is assumed, so the screen applies a fortiori to
  fully-invested long-only books.

  \Cref{chap:distance-implied} carries its own label for this declaration under
  the distance-implied reading; the two manuscripts consume one theorem.
\end{theorem}

\begin{proof}
  \uses{cor:approximate-map-bound}
  Sum the pairwise floor against the nonnegative weight products.
\end{proof}
